% BORRADOR para revisar: no está incluido en main.tex. Si Godoy lo prefiere dentro de cap:prototipo, basta
% bajar un nivel cada título. Cifras medidas en el repositorio el 2026-11-04. Si se aprueba, conviene llevar
% aquí el último párrafo de sec:prototipo-software, para no repetirlo.
\chapter{Desarrollo y despliegue del sistema}
\label{cap:desarrollo}

El Capítulo~\ref{cap:prototipo} describe qué hace el sistema y por qué es seguro. Este capítulo describe
cómo se construyó: los requisitos, la arquitectura de software, los problemas de implementación, la
verificación y la forma de instalarlo.

\section{Proceso de desarrollo}
\label{sec:desarrollo-proceso}

El desarrollo fue iterativo y se registró en un repositorio git, con 77 \textit{commits} entre el 19 de
septiembre y el 4 de noviembre de 2026. El diseño cambió de forma importante en ese período. La primera versión
controlaba el acceso a un recinto residencial; después se adaptó a eventos masivos. En una versión temprana el
servidor emitía un código opaco, y el 29 de septiembre el QR pasó a generarse en el teléfono, con la cédula
firmada por el emisor y cifrado para los lectores autorizados.

Cada decisión de diseño quedó en un registro de 18 decisiones (\textit{Architecture Decision Records}), que
conserva también las reemplazadas y sus motivos. El comportamiento esperado de cada módulo está en
especificaciones con requisitos y escenarios, escritas con OpenSpec, a partir de las cuales se construye la
Tabla~\ref{tab:desarrollo-rf}.

El desarrollo se apoyó en el asistente de programación Claude Code (Anthropic): el historial del repositorio
registra coautoría del asistente en 68 de los 77 \textit{commits}.
% TODO (autor): DECISIÓN IMPORTANTE. sec:prototipo-software solo declara el uso del asistente en las revisiones
% de código, pero el repositorio es público y muestra esta coautoría. Revisa la política de la universidad,
% describe con tus palabras qué hiciste tú y qué hizo el asistente, y deja ambas secciones coherentes. No
% atribuí la autoría de las decisiones: eso solo lo puedes afirmar tú.

\section{Requisitos}
\label{sec:desarrollo-requisitos}

Los requisitos funcionales salen del caso de uso y del modelo de amenazas. La
Tabla~\ref{tab:desarrollo-rf} los resume junto con las pruebas que los verifican, y la
Tabla~\ref{tab:desarrollo-rnf} hace lo mismo con los no funcionales.

\begin{table}[htbp]
  \centering
  \caption{Requisitos funcionales y dónde se verifican.}
  \label{tab:desarrollo-rf}
  \small
  \begin{tabular}{lp{7.6cm}p{5.0cm}}
    \toprule
    Id & Requisito & Pruebas \\
    \midrule
    RF1 & Activar la cédula con RUT y clave, ligada a un teléfono & \texttt{test\_registro\_civil} \\
    RF2 & Guardar la cédula cifrada con un PIN, con bloqueo por intentos & \texttt{wallet.test} \\
    RF3 & Generar el QR en el teléfono, sin conexión, cada 45~s & \texttt{credential.test}, \texttt{CedulaPage.test} \\
    RF4 & Verificar firmas y mostrar la fotografía en el lector & \texttt{test\_qr\_flow}, \texttt{credential.test} \\
    RF5 & Verificar en dos pasos, con la decisión del guardia & \texttt{test\_qr\_flow}, \texttt{ScannerPage.test} \\
    RF6 & Impedir el reúso de un QR y un segundo ingreso & \texttt{test\_qr\_flow} \\
    RF7 & Administrar entradas nominativas y prohibiciones & \texttt{test\_tickets}, \texttt{test\_bans} \\
    RF8 & Revocar cédulas con una lista firmada & \texttt{test\_revocation} \\
    RF9 & Operar sin red y sincronizar, con alertas & \texttt{test\_offline}, \texttt{offline.test} \\
    RF10 & Registrar el historial y las métricas por fila & \texttt{test\_logs} \\
    RF11 & Administrar al personal por roles & \texttt{test\_users} \\
    RF12 & Publicar lo que un verificador de terceros necesita & \texttt{test\_verifier\_spec} \\
    RF13 & Ofrecer la comparación facial opcional, con consentimiento & \texttt{test\_informe\_gaps}, \texttt{ScannerPage.test} \\
    \bottomrule
  \end{tabular}
\end{table}

\begin{table}[htbp]
  \centering
  \caption{Requisitos no funcionales y cómo se comprueban.}
  \label{tab:desarrollo-rnf}
  \small
  \begin{tabular}{p{3.2cm}p{6.0cm}p{4.2cm}}
    \toprule
    Atributo & Requisito & Comprobación \\
    \midrule
    Seguridad de la credencial & Una copia diferida del QR no permite entrar & Escenarios de ataque (Sección~\ref{sec:prototipo-ataques}) \\
    Seguridad de la aplicación & Controles según OWASP Top~10~\cite{owasp_top10_2021} & \texttt{test\_security} \\
    Privacidad & El recinto no recibe la fotografía; el historial enmascara el RUT & \texttt{test\_qr\_flow}, \texttt{test\_logs} \\
    Concurrencia & Ninguna doble concesión con varias filas a la vez & Prueba de carga (Sección~\ref{sec:prototipo-fila}) \\
    Disponibilidad & La puerta sigue operando sin red & \texttt{test\_offline} \\
    Compatibilidad & Teléfono y lector producen y leen el mismo QR & Vectores compartidos \\
    Usabilidad & La interfaz recoge los arquetipos de la encuesta & \texttt{CedulaPage.test} (Tabla~\ref{tab:trazabilidad-anexo}) \\
    \bottomrule
  \end{tabular}
\end{table}

La usabilidad con personas todavía no se ha medido. El piloto entrega la primera evidencia.
% TODO (autor): ajustar esta frase cuando exista la sección del piloto.

\section{Arquitectura de software}
\label{sec:desarrollo-arquitectura}

El sistema tiene tres componentes que se ejecutan como procesos separados, con los actores de la
Figura~\ref{fig:prototipo-arquitectura}. El \emph{Registro Civil simulado} es un servicio aparte, con unas pocas
rutas, que guarda la clave privada del emisor y el padrón. El \emph{backend del recinto} guarda las entradas, las
prohibiciones y el historial, y solo conoce la clave pública del emisor. El \emph{frontend} es una sola aplicación
web que contiene dos experiencias: la cédula del asistente, en una ruta pública, y el escáner y la administración
del personal, en rutas protegidas por rol. Como la criptografía del teléfono y del lector corre en el navegador,
no hay que instalar una aplicación nativa. La Tabla~\ref{tab:desarrollo-tecnologias} resume las tecnologías.

\begin{table}[htbp]
  \centering
  \caption{Tecnologías usadas.}
  \label{tab:desarrollo-tecnologias}
  \small
  \begin{tabular}{p{3.0cm}p{5.0cm}p{5.4cm}}
    \toprule
    Capa & Tecnología & Uso \\
    \midrule
    Servicios & Python, FastAPI, Pydantic & API y validación de cada entrada \\
    Datos & SQLAlchemy, SQLite & Persistencia; SQLite se declara no apto para producción a escala \\
    Criptografía del servidor & \texttt{cryptography} & Ed25519, X25519, HKDF y ChaCha20-Poly1305 \\
    Sesiones & Argon2id (\texttt{argon2-cffi}), PyJWT & Contraseñas y tokens del personal \\
    Interfaz & React, TypeScript, Vite, Tailwind CSS & Aplicación web \\
    Criptografía del cliente & \texttt{@noble/*} & Las mismas primitivas, en el navegador \\
    Lectura del QR & \texttt{html5-qrcode}, \texttt{barcode-detector} & Cámara e imágenes, con ZXing-C++ en WebAssembly \\
    Comparación facial & \texttt{onnxruntime-web} & Módulo opcional, en el dispositivo del guardia \\
    Pruebas y calidad & \texttt{pytest}, Vitest, oxlint & Pruebas automáticas y análisis estático \\
    \bottomrule
  \end{tabular}
\end{table}
% TODO (autor): si quieres justificar por qué elegiste FastAPI y React frente a otras opciones, agrégalo aquí. No
% está documentado en el repositorio, así que no lo inventé.

\subsection{Modelo de datos}
\label{subsec:desarrollo-datos}

La Tabla~\ref{tab:desarrollo-datos} resume las tablas del recinto. La fotografía no aparece en ninguna, y tampoco
hay imágenes ni vectores faciales. Las tablas \texttt{qr\_uses} y \texttt{qr\_checks} usan el \textit{nonce} como
clave primaria, y de eso depende la atomicidad del uso único.
El padrón del emisor queda fuera de esta base, en un archivo del servicio simulado.

\begin{table}[htbp]
  \centering
  \caption{Tablas del sistema del recinto.}
  \label{tab:desarrollo-datos}
  \small
  \begin{tabular}{p{3.6cm}p{10.0cm}}
    \toprule
    Tabla & Contenido \\
    \midrule
    \texttt{users} & Personal: rol, contraseña y versión de sesión \\
    \texttt{events}, \texttt{gates} & Eventos y sus filas, digital o normal \\
    \texttt{tickets} & Entrada nominativa: RUT, evento y marca de ingreso \\
    \texttt{banned\_persons} & Prohibiciones de ingreso, por RUT \\
    \texttt{qr\_uses} & \textit{Nonces} ya usados \\
    \texttt{qr\_checks} & QR verificados, a la espera de la decisión del guardia \\
    \texttt{offline\_syncs} & Decisiones sin red ya sincronizadas \\
    \texttt{biometric\_consents} & Evidencia del consentimiento: \textit{nonce}, fila, quién y cuándo \\
    \texttt{access\_logs} & Historial de cada intento, con motivo, latencia y marca sin conexión \\
    \bottomrule
  \end{tabular}
\end{table}

\subsection{Interfaces}
\label{subsec:desarrollo-api}

El backend del recinto expone 24 rutas HTTP, agrupadas en la Tabla~\ref{tab:desarrollo-api}, y el emisor simulado
expone 7. El emisor publica su clave pública y su lista de revocadas sin autenticación; el formato del QR está en
la especificación para verificadores del repositorio, y un verificador escrito solo a partir de ese documento
pasa las pruebas. Ese verificador necesita además la clave de lectores autorizados, que entrega el emisor.
% TODO (autor): coordinar esta última frase con la corrección de conclu.tex (pregunta 6, aporte y cuarta
% recomendación), que hoy enumera solo tres elementos.

\begin{table}[htbp]
  \centering
  \caption{Rutas del backend del recinto.}
  \label{tab:desarrollo-api}
  \small
  \begin{tabular}{p{3.2cm}cp{4.4cm}p{4.4cm}}
    \toprule
    Grupo & Rutas & Acceso & Función \\
    \midrule
    \texttt{/auth} & 2 & Público y personal & Inicio de sesión \\
    \texttt{/users} & 4 & Administración & Personal \\
    \texttt{/events} & 7 & Personal; una pública & Eventos, filas, claves y copia sin red \\
    Entradas & 3 & Administración & Lista de RUT por evento \\
    \texttt{/qr} & 2 & Personal & Validación y sincronización \\
    \texttt{/bans} & 3 & Administración & Prohibiciones \\
    \texttt{/logs}, \texttt{/metrics} & 2 & Personal & Historial y métricas \\
    \texttt{/health} & 1 & Público & Estado del servicio \\
    \bottomrule
  \end{tabular}
\end{table}

\section{Problemas de implementación}
\label{sec:desarrollo-implementacion}

Varias decisiones salieron de problemas encontrados al construir y probar el sistema.

\emph{Mismo formato en dos lenguajes.} El QR se arma en TypeScript en el teléfono, y el backend lo verifica en
Python. Cualquier diferencia de un byte invalidaría las firmas, de modo que ambos lados se prueban contra un mismo
archivo de vectores y deben producir el mismo QR byte a byte.

\emph{Lectura de un QR denso.} Con la fotografía, el QR queda cerca de la versión~25, y el decodificador de
respaldo de \texttt{html5-qrcode} no lo leía ni desde una imagen limpia. Se resolvió con la interfaz
\texttt{BarcodeDetector}, usando ZXing-C++ en WebAssembly donde el navegador no la ofrece.

\emph{Pistolas lectoras.} Una pistola 2D en modo teclado escribe el QR como si fueran teclas. Si emula un teclado
estadounidense en un Windows en español, los símbolos llegan cambiados, así que el escáner detecta el prefijo
alterado y lo corrige. Las pistolas que no emulan teclado se leen por puerto serie con Web Serial.

\emph{Concurrencia en el uso único.} Con varias filas en paralelo, consultar y después insertar permite dos
ingresos con el mismo QR. El uso único se resolvió insertando el \textit{nonce} como clave primaria y decidiendo
según si la inserción tuvo efecto. La sincronización sin red usa la misma técnica, porque dos envíos simultáneos
registraban el ingreso dos veces y generaban una alerta falsa.

\emph{Hora declarada por el lector.} Sin red, el escáner informa la hora de cada escaneo, y una versión anterior
aceptaba cualquier hora pasada. El servidor la acota ahora con la copia sellada de contingencia
(Sección~\ref{sec:prototipo-contingencia}).

\section{Verificación del software}
\label{sec:desarrollo-verificacion}

El software tiene 271 pruebas automáticas, 162 del servidor y 109 del cliente.
Hay pruebas unitarias, de integración contra la API con una base temporal y de interfaz con la aplicación
renderizada, además de los vectores compartidos.
% TODO (autor): actualizar estas cifras al cerrar el código. sec:prototipo-software dice 235, que corresponde a
% 77c29fa. Si esta sección se mantiene, unificar y citar el commit final.
La prueba de la especificación para terceros reimplementa el verificador sin importar el código del proyecto, de
modo que falla si el formato cambia sin actualizar el documento.

Las pruebas y el análisis estático se ejecutan a mano; no hay integración continua. El 9 de octubre de 2026 una
revisión completa del código encontró 34 problemas, y una segunda revisión, de los arreglos, otros 8
(Sección~\ref{sec:prototipo-software}). La evaluación de seguridad y de carga se describe en la
Sección~\ref{sec:prototipo-metodo}.

\section{Instalación y despliegue}
\label{sec:desarrollo-despliegue}

El sistema se ha ejecutado en el equipo del autor, en los tres modos de la
Tabla~\ref{tab:desarrollo-modos}, que el mismo programa de instalación levanta con un comando. Ese programa instala
las dependencias, crea un secreto aleatorio para las sesiones si falta, prepara los datos de demostración y
compila el cliente. Las claves del emisor quedan fuera de git. El servidor no arranca sin un secreto de al menos 32
caracteres.

\begin{table}[htbp]
  \centering
  \caption{Modos de ejecución del sistema.}
  \label{tab:desarrollo-modos}
  \small
  \begin{tabular}{p{2.6cm}p{4.0cm}p{6.8cm}}
    \toprule
    Modo & Alcance & Uso \\
    \midrule
    Local & Solo el equipo & Desarrollo y pruebas \\
    Red local & Equipos de la misma red, con HTTPS autofirmado y sin internet & Piloto en terreno \\
    Túnel temporal & Público, con una dirección que cambia en cada ejecución & Demostraciones breves \\
    \bottomrule
  \end{tabular}
\end{table}

El modo de red local requiere HTTPS porque los navegadores móviles solo dan acceso a la cámara en un contexto
seguro. Cada teléfono debe aceptar la advertencia del certificado autofirmado, un costo aceptable en un piloto pero
no en un evento.

El sistema \emph{no se ha desplegado} en un servidor permanente, decisión que se pospuso hasta la aprobación del
diseño. Un despliegue real exigiría, como mínimo, TLS con un certificado válido, una base de datos con varios
escritores en lugar de SQLite, un almacén seguro para las claves, límites de intentos compartidos entre procesos,
un registro de seguridad centralizado, un plazo de conservación de los datos personales y un segundo factor para el
personal. La mayoría ya se declaró como limitación en la Sección~\ref{sec:prototipo-limitaciones}. Un emisor real
reemplazaría además al servicio simulado, lo que hoy depende de que el Registro Civil publique su interfaz.
% TODO (autor): si después de la aprobación se despliega, describir aquí la infraestructura real.
